\scanpage{63}
\diagram{\begin{tikzcd}
X \arrow[r,phantom,"\subset" description] \arrow[d] & Y \arrow[dl] \\
S &
\end{tikzcd}}
avec $Y$ form\textsuperscript{t}. lisse sur $S$ [tjs possible \deleted{loc\textsuperscript{t}. sur $X$}, \unsure{ou} si $S,X$ affines, ou si $X$ est qu. projectif sur $S$ \deleted{\unsure{de t. f.}} admettant \unsure{f. inv.} ample]. Considérons le foncteur sur le site cristallin de $X/S$ ``représenté'' par $(X,Y)$ ; de façon précise, si $Y(i)$ est le $i$ième vois. inf. de $X$ dans $Y$, on trouve que ce foncteur est le faisceau
\[
\widetilde Y=\varinjlim_i\widetilde{Y(i)},
\]
d'autre part l'hyp que $Y$ est \unsure{f\textsuperscript{t}} lisse sur $S$ implique que $\widetilde Y$ couvre le faisceau final. D'où encore
\[
H^*(X_{\mathrm{cris}},F)\Longleftarrow E_2^{pq}=H^p(\nu\longmapsto H^q(\widetilde Y^{\nu+1},F))
\]
\mspage{62}
